\documentclass[10pt,a4paper]{amsart}
\usepackage[margin=19mm]{geometry}
\usepackage{amsmath,amssymb,mathtools}
\usepackage[scr=rsfs,cal=euler]{mathalfa}
\usepackage{xfrac}
\usepackage[unicode,hidelinks,bookmarks=false]{hyperref}
\newcommand{\ie}{i.e.\ }
\newcommand{\cf}{cf.\ }
\newcommand{\resp}{respectively\ }
\newcommand{\Subsection}{\subsection}
\providecommand{\nobreakdash}{\nobreak\hbox{-}}
\setlength{\emergencystretch}{2em}
\multlinegap0pt
\usepackage{bm}
\usepackage{bbm}
\usepackage{dsfont}
\usepackage{xspace}
\usepackage{cancel}
\usepackage{mathtools}
\usepackage[shortlabels]{enumitem}
\usepackage{amsthm}
\makeatletter
\@ifundefined{thm}{\newtheorem{thm}{Thm}}{}
\@ifundefined{lemma}{\newtheorem{lemma}[thm]{Lemma}}{}
\makeatother
\DeclareMathOperator{\Cov}{Cov}
\DeclareMathOperator{\EE}{\mathbb{E}}
\DeclareMathOperator{\QQ}{\mathbb{Q}}
\DeclareMathOperator{\Var}{Var}
\DeclareMathOperator{\sumij}{\sum_{\mathnormal{ i,j}=1}^{\mathnormal n}}
\DeclareMathOperator{\sumin}{\sum_{\mathnormal i=1}^{\mathnormal n}}
\DeclareMathOperator{\sumjn}{\sum_{\mathnormal j=1}^{\mathnormal n}}
\newcommand{\A}{\mathbf{A}}
\newcommand{\X}{\mathbf{X}}
\newcommand{\Z}{\mathbf{Z}}
\newcommand{\an}{\alpha_{\mathnormal n}}
\newcommand{\ttheta}{\boldsymbol{\theta}}
\title[Inference on Gaussian mixture models with dependent labels]{Inference on Gaussian mixture models with dependent labels}
\author{Seunghyun Lee, Rajarshi Mukherjee, Sumit Mukherjee}
\date{Source: arXiv:2510.06501v1 (CC BY 4.0)}
\begin{document}
\maketitle

\noindent\emph{Source and licence.} Inference on Gaussian mixture models with dependent labels. Version arXiv:2510.06501v1 (CC BY 4.0), distributed under the Creative Commons Attribution 4.0 International licence, as recorded in the arXiv licence metadata for this exact version and shown on the abstract page https://arxiv.org/abs/2510.06501v1. Full rights notice: licenses/arxiv-2510.06501-CC-BY-4.0.txt.\par\smallskip

\noindent\textit{Editorial scope.} Counted unit: the Lemma whose source label is \texttt{lem:Ising moment concentration} in the article (source0.tex lines 2172--2181, source label \texttt{lem:Ising moment concentration}), delivered together with the article's own complete original proof of it (source0.tex lines 2182--2219, one \texttt{proof} environment of 38 lines). No mathematics is written, repaired or re-derived: statement and proof are the article's own text line for line. Required context retained uncounted: eq:mean field assumption (lines 485--487); lem:conditional sum (lines 1043--1054). Every definition, display and label of those lines is retained unchanged. Omitted from this package and not redistributed: 1--484, 488--1042, 1055--2171, 2220--2739. Nothing inside a retained excerpt is omitted or altered: every retained body is delivered exactly as the article's lines are written, so no omission receipt of any kind accompanies this package. Citations: the retained excerpts carry no citation command at all, so no bibliography entry of the article is transcribed and no reference is redistributed. Reference adaptation: none was needed and none was made; every reference inside the delivered unit resolves inside it, so no unresolved reference or citation is produced. Printed slips kept literal: no printed word, symbol, hypothesis or number of the counted statement or of its proof was altered. Layout and class: the article's own class is \texttt{imsart} with packages including amsthm, amsmath, amsfonts, amssymb, natbib, graphicx, enumitem, bbm, comment, hyperref, subcaption, cleveref, booktabs, multirow; the wrapper is the lane's pinned \texttt{amsart} editorial wrapper at 10pt on A4, which restates the theorem-like environments the retained text needs and adds the article's own macro definitions that the retained text needs (\texttt{\textbackslash A}, \texttt{\textbackslash Cov}, \texttt{\textbackslash EE}, \texttt{\textbackslash QQ}, \texttt{\textbackslash Var}, \texttt{\textbackslash X}, \texttt{\textbackslash Z}, \texttt{\textbackslash an}, \texttt{\textbackslash sumij}, \texttt{\textbackslash sumin}, \texttt{\textbackslash sumjn}, \texttt{\textbackslash ttheta}), with extra wrapper packages (bm, bbm, dsfont, xspace, cancel, mathtools, enumitem). The delivered numbering is the unit's own. Assets: no retained span contains a figure environment, a graphics-inclusion command, a PDF-page inclusion, a table or a TikZ picture; no image file is copied into this package, the package declares no asset dependency and no image file is redistributed with it. Licence: version arXiv:2510.06501v1 of this article is distributed under the Creative Commons Attribution 4.0 International licence, as recorded in the arXiv licence metadata for this exact version and shown on the abstract page https://arxiv.org/abs/2510.06501v1. A full rights notice is delivered at licenses/arxiv-2510.06501-CC-BY-4.0.txt. Characters: every retained excerpt is pure ASCII, so the delivered engine, XeLaTeX with its default Latin Modern text font, reports no missing character.
    \begin{equation}\label{eq:mean field assumption}
    \alpha_n := \max_{i=1}^n \sum_{j=1}^n A_n(i,j)^2 =o \left( \frac{1}{\sqrt{n \log n}} \right).
    \end{equation}

\begin{lemma}\label{lem:conditional sum}
    Suppose $\beta <1$, $\Z^n \sim \QQ_{0,\beta, \A_n}$, and $\X^n \sim P_{\ttheta_0,\QQ_{0,\beta, \A_n}}$. Then, the following holds, where $\an = \max_{i=1}^n \sumjn A_n(i,j)^2$ and the hidden constants only depend on $K, C$.
    \begin{enumerate}[(a)]
    \item Let $\phi: \mathbb{R}^d \to \mathbb{R}$ be an odd function with $|\EE(\phi(\X) | Z = z)| \le K$ and $\Var(\phi(\X) | Z = z) \le C$ for $K, C < \infty$ and $z = \pm 1$. Then,
    \begin{equation*}
        \EE \sumin \left(\sumjn A_n(i,j) \phi(\X_j) \right)^2 = O(n \an).
        %
    \end{equation*}
    \item Let $\phi_1, \phi_2: \mathbb{R}^d \to \mathbb{R}$ be odd functions with $|\EE(\phi_a(\X) | Z = z)| \le K_a$ and $\Var(\phi_a(\X) | Z = z) \le C$, $\Cov (\phi_1(\X), \phi_2(\X) | Z = z) \le C$ for $a = 1,2$, and $z = \pm 1$, where $K_a, C < \infty$. Then,
    $$\EE \left( \sumij A_n(i,j) \phi_1(\X_i) \phi_2(\X_j) \right)^2 = O(n^2 \alpha_n^2 + n\an).$$
    \end{enumerate}
    \end{lemma}

\begin{lemma}\label{lem:Ising moment concentration}
    Suppose $\beta < 1$ and let $\Z^n \sim \QQ_{0, \beta, \A_n}$. Then, the following bounds hold.
    \begin{enumerate}[(a)]
        \item $\EE (\Z^\top \A_n \Z)^2 = O(n^2 \alpha_n^2 + n\an)$
        \item $\EE (\Z^\top \A_n^2 \Z)^2 %
        = O(n^2 \alpha_n^2)$.
        %
    \end{enumerate}
    In particular, under \eqref{eq:mean field assumption}, the RHS of (a) and (b) can be replaced with $o(n)$. %
\end{lemma}

\begin{proof}[Proof of Lemma \ref{lem:conditional sum}]
    \begin{enumerate}[(a)]
    \item Define $m_i({\Z^n}) := \sumjn A_n(i,j) Z_j$
    and note that
    \begin{equation}
        |\sumjn A_n(i,j) \phi(\X_j)| \le |\sumjn A_n(i,j) (\phi(\X_j) - K Z_j)| + |K| |m_i({\Z^n})|.
    \end{equation}
    Since $\X^n \mid \Z^n$ is independent, we have
    \begin{align*}
        \EE \left[ |\sumjn A_n(i,j) (\phi(\X_j) - K Z_j)|^2 \mid \Z^n \right] & = \Var \left( \sumjn A_n(i,j) \phi(\X_j) \mid \Z^n \right) \\
        = \sumjn A_n(i,j)^2 &\Var( \phi(\X_j) \mid \Z^n ) \le C \sumjn A_n(i,j)^2
    \end{align*}
    for all $i$. Also, Lemma \ref{lem:Ising moment concentration}(b) gives $\EE (\Z^\top \A_n^2 \Z)^2 = \EE (\sumin m_i^2(\Z))^2 \lesssim n\an$.
    The proof is complete by summing the two bounds.
   
    \item By expanding the square and using the independence of $\X^n \mid \Z^n$, we have
    \begin{equation}\label{eq:2nd moment conditional expectation}
    \begin{aligned}
        &\EE \left[ \left( \sumij A_n(i,j) \phi_1(\X_i) \phi_2(\X_j) \right)^2 \mid \Z^n \right] \\
        =&  \EE \left[ \sum_{i,j,k,l} A_n(i,j) A_n(k,l) \phi_1(\X_i) \phi_2(\X_j) \phi_1(\X_k) \phi_2(\X_l) \mid \Z^n \right]  \\
        \lesssim& |\sum_{|\{i,j,k,l\}|=4} A_n(i,j) A_n(k,l) Z_i Z_j Z_k Z_l| + |\sum_{i=k \neq j \neq l} A_n(i,j) A_n(i,l) Z_i^2 Z_j Z_l| \\
        & \quad + |\sum_{i, j} A_n(i,j)^2 Z_i^2 Z_j^2|.
    \end{aligned}
    \end{equation}
    Here, we have omitted displaying the constants arising from moments of $\phi_1, \phi_2$, which only depend on $K, C$.
    Noting that $|Z_i| = 1$, it is easy to control the last two terms:
    \begin{align*}
        \sum_{i=k \neq j \neq l} A_n(i,j) A_n(i,l) Z_i^2 Z_j Z_l &= \sum_{i=k \neq j \neq l} A_n(i,j) A_n(i,l) Z_j Z_l = \Z^\top \A_n^2 \Z, \\
        \sum_{i, j} A_n(i,j)^2 Z_i^2 Z_j^2 &= \sum_{i, j} A_n(i,j)^2 \lesssim n\an.
    \end{align*}
    Hence, by taking an expectation over $\Z^n$ on \eqref{eq:2nd moment conditional expectation} and using Lemma \ref{lem:Ising moment concentration},
    \begin{align*}
        \EE \left( \sumij A_n(i,j) \phi_1(\X_i) \phi_2(\X_j) \right)^2 &\le \EE |\sum_{|\{i,j,k,l\}|=4} A_n(i,j) A_n(k,l) Z_i Z_j Z_k Z_l| + O(n\an) \\
        &\le \EE |\sum_{i,j,k,l} A_n(i,j) A_n(k,l) Z_i Z_j Z_k Z_l| + O(n\an) \\
        &= \EE \left(\Z^\top \A_n \Z \right)^2 + O(n\an) = O(n^2 \alpha_n^2 + n\an).
    \end{align*}
    \end{enumerate}
    \end{proof} %

\end{document}
